% Chapter of the map: Bounded arithmetic and logic.
% Conventions and the template of an entry: blueprint/README.md. Check with
% blueprint/tools/check_names.sh. Labels of this file begin with kind:arithmetic (def:arithmetic-..., thm:arithmetic-...).
\chapter{Bounded arithmetic and logic}
\label{chap:arithmetic}

This chapter has the formalization of Buss's bounded arithmetic $S^1_2$ and of the results around it.
The language, the axioms BASIC, the $\Sigma^b_1$ formulas and the theory $S^1_2$ are written on the
first-order logic of Mathlib, and so are $T^1_2$ and $T^0_2$. A sequent calculus with PIND as a rule
and the completeness theorem for it come next. Then come the Cobham functions, the witnessing theorem
and the theorem of Parikh, with the results about polynomial time that follow from them. The last
sections are about the statements that $S^1_2$ and $T^1_2$ do not prove if some conjecture holds,
about a propositional version of $S^1_2$ in the system of Frege, and about proof systems. Many results
are stated in a model smaller than the one of the source (a sequent calculus, the theory of Cobham
programs, an abstract proof system); each entry says which.

\section{The theory $S^1_2$}
\label{sec:arithmetic-s12}

The theory $S^1_2$ is given as a class of structures of the language of bounded arithmetic, and also
as a theory of Mathlib. The two are shown to have the same models.

\begin{definition}[The language of bounded arithmetic and the axioms BASIC]
  \label{def:arithmetic-basic}
  \lean{Sagitta.Map.Acotada.L, Sagitta.Map.Acotada.EsModeloBASIC}
  \leanok
  \own{}
  \infrastructure{}
  The language $L$ has the constant $0$, the unary function symbols $S$, $|x|$ and $\lfloor x/2
  \rfloor$, the binary function symbols $+$, $\cdot$ and $x \# y$, and the binary relation symbol
  $\le$. Write $1$ for $S0$ and $2x$ for $(SS0)\cdot x$. An $L$-structure satisfies BASIC when, for
  all its elements $x, y, z, u, v$, it satisfies the following $32$ clauses.
  \begin{enumerate}
    \item $y \le x \to y \le Sx$.
    \item $x \ne Sx$.
    \item $0 \le x$.
    \item $(x \le y \wedge x \ne y) \leftrightarrow Sx \le y$.
    \item $x \ne 0 \to 2x \ne 0$.
    \item $y \le x \vee x \le y$.
    \item $x \le y \to y \le x \to x = y$.
    \item $x \le y \to y \le z \to x \le z$.
    \item $|0| = 0$.
    \item $x \ne 0 \to |2x| = S|x| \wedge |S(2x)| = S|x|$.
    \item $|1| = 1$.
    \item $x \le y \to |x| \le |y|$.
    \item $|x \# y| = S(|x| \cdot |y|)$.
    \item $0 \# y = 1$.
    \item $x \ne 0 \to 1 \# (2x) = 2(1 \# x) \wedge 1 \# S(2x) = 2(1 \# x)$.
    \item $x \# y = y \# x$.
    \item $|x| = |y| \to x \# z = y \# z$.
    \item $|x| = |u| + |v| \to x \# y = (u \# y)\cdot(v \# y)$.
    \item $x \le x + y$.
    \item $x \le y \to x \ne y \to S(2x) \le 2y \wedge S(2x) \ne 2y$.
    \item $x + y = y + x$.
    \item $x + 0 = x$.
    \item $x + Sy = S(x + y)$.
    \item $(x + y) + z = x + (y + z)$.
    \item $x + y \le x + z \leftrightarrow y \le z$.
    \item $x \cdot 0 = 0$.
    \item $x \cdot Sy = x \cdot y + x$.
    \item $x \cdot y = y \cdot x$.
    \item $x \cdot (y + z) = x \cdot y + x \cdot z$.
    \item $1 \le x \to (x \cdot y \le x \cdot z \leftrightarrow y \le z)$.
    \item $x \ne 0 \to |x| = S|\lfloor x/2 \rfloor|$.
    \item $x = \lfloor y/2 \rfloor \leftrightarrow (2x = y \vee S(2x) = y)$.
  \end{enumerate}
  \leanfile{Sagitta/Map/Acotada.lean}
\end{definition}

\begin{definition}[Sharply bounded, $\Sigma^b_1$ and $\Pi^b_1$ formulas]
  \label{def:arithmetic-sigma1}
  \lean{Sagitta.Map.Acotada.EsAguda, Sagitta.Map.Acotada.EsSigma1, Sagitta.Map.Acotada.EsPi1}
  \leanok
  \uses{def:arithmetic-basic}
  \own{}
  \infrastructure{}
  Write $\forall x \le t\, \varphi$ for $\forall x (x \le t \to \varphi)$ and $\exists x \le t\,
  \varphi$ for $\exists x (x \le t \wedge \varphi)$. The \emph{sharply bounded} formulas of $L$ are
  the smallest class that contains $\bot$ and the atomic formulas, is closed under implication, and
  is closed under $\forall x \le |t|$ and $\exists x \le |t|$ for every term $t$. The classes of
  $\Sigma^b_1$ and $\Pi^b_1$ formulas are defined together. The $\Sigma^b_1$ formulas are the
  smallest class such that: (a) it contains the sharply bounded formulas; (b) it is closed under
  $\exists x \le t$ and under $\forall x \le |t|$, for every term $t$; (c) it is closed under
  $\wedge$ and $\vee$; (d) if $B$ is $\Pi^b_1$ and $A$ is $\Sigma^b_1$, then $\neg B$ and $B \to A$
  are $\Sigma^b_1$. The $\Pi^b_1$ formulas are the smallest class such that: (a) it contains the
  sharply bounded formulas; (b) it is closed under $\forall x \le t$ and under $\exists x \le |t|$;
  (c) it is closed under $\wedge$ and $\vee$; (d) if $B$ is $\Sigma^b_1$ and $A$ is $\Pi^b_1$, then
  $\neg B$ and $B \to A$ are $\Pi^b_1$.
  \leanfile{Sagitta/Map/Acotada.lean}
\end{definition}

\begin{definition}[Models of $S^1_2$]
  \label{def:arithmetic-s12-model}
  \lean{Sagitta.Map.Acotada.EsModeloS12}
  \leanok
  \uses{def:arithmetic-basic, def:arithmetic-sigma1}
  \source{\cite{Buss86}.}
  \fidelity{Not checked against the printed source.}
  An $L$-structure $M$ is a model of $S^1_2$ when it satisfies BASIC and the induction on the
  length (PIND) for $\Sigma^b_1$ formulas: for every $k$, every $\Sigma^b_1$ formula $\varphi(x,
  p_1, \dots, p_k)$ and all $p \in M^k$, if $\varphi(0, p)$ holds and $\varphi(\lfloor x/2 \rfloor,
  p) \to \varphi(x, p)$ holds for every $x \in M$, then $\varphi(x, p)$ holds for every $x \in M$.
  \leanfile{Sagitta/Map/Acotada.lean}
\end{definition}

\begin{theorem}[The standard structure is a model of $S^1_2$]
  \label{thm:arithmetic-nat-s12}
  \lean{Sagitta.Map.Acotada.nat_esModeloS12}
  \leanok
  \uses{def:arithmetic-s12-model}
  Let $\mathbb{N}$ be the $L$-structure in which $|x|$ is the number of binary digits of $x$, $x \#
  y = 2^{|x| \cdot |y|}$, $\lfloor x/2 \rfloor$ is the integer part of $x/2$, and the other symbols
  have their usual meaning. Then $\mathbb{N}$ is a model of $S^1_2$.
  \source{\cite{Buss86}.}
  \fidelity{Not checked against the printed source.}
  \leanfile{Sagitta/Map/Acotada.lean}
\end{theorem}

\begin{proof}
  \leanok
  Each of the $32$ clauses of BASIC is checked in $\mathbb{N}$ by simplification and \texttt{omega}.
  PIND holds for every formula, by strong induction on $x$, because $\lfloor x/2 \rfloor < x$ when $x
  \ne 0$.
\end{proof}

\begin{theorem}[$S^1_2$ as a theory of Mathlib has the same models]
  \label{thm:arithmetic-s12-theory}
  \lean{Sagitta.Map.Acotada.modelo_S12_iff}
  \leanok
  \uses{def:arithmetic-s12-model}
  Let $S^1_2$ also denote the theory of $L$ made of the $32$ axioms of BASIC and, for each
  $\Sigma^b_1$ formula $\varphi(x, p_1, \dots, p_k)$, the universal closure of the PIND axiom for
  $\varphi$. For every $L$-structure $M$, $M$ satisfies this theory if and only if $M$ is a model of
  $S^1_2$ in the sense of Definition~\ref{def:arithmetic-s12-model}.
  \source{\cite{Buss86}.}
  \fidelity{Not checked against the printed source.}
  \leanfile{Sagitta/Map/Teoria.lean}
\end{theorem}

\begin{proof}
  \leanok
  The theory is a union of the axioms of BASIC and the PIND axioms. The BASIC part is the
  equivalence for BASIC, and for each PIND axiom, realizing its closure is the clause of the
  definition.
\end{proof}

\section{$T^1_2$, $T^0_2$ and their inclusions with $S^1_2$}
\label{sec:arithmetic-t12}

The theories $T^1_2$ and $T^0_2$ are given by their models, in the same style as $S^1_2$. The two
inclusions are statements about all models.

\begin{definition}[Models of $T^1_2$ and $T^0_2$]
  \label{def:arithmetic-t12-models}
  \lean{Sagitta.Map.Acotada.EsModeloT12, Sagitta.Map.Acotada.EsModeloT02}
  \leanok
  \uses{def:arithmetic-basic, def:arithmetic-sigma1}
  \own{}
  \infrastructure{}
  An $L$-structure $M$ is a model of $T^1_2$ when it satisfies BASIC and the ordinary induction for
  $\Sigma^b_1$ formulas: for every $k$, every $\Sigma^b_1$ formula $\varphi(x, p_1, \dots, p_k)$ and
  all $p \in M^k$, if $\varphi(0, p)$ holds and $\varphi(x, p) \to \varphi(Sx, p)$ holds for every
  $x \in M$, then $\varphi(x, p)$ holds for every $x \in M$. It is a model of $T^0_2$ when it
  satisfies BASIC and the same induction for the sharply bounded formulas.
  \leanfile{Sagitta/Map/T12.lean}
\end{definition}

\begin{theorem}[Every model of $S^1_2$ is a model of $T^0_2$]
  \label{thm:arithmetic-s12-contains-t02}
  \lean{Sagitta.Map.Acotada.Arranque.s12ContieneT02}
  \leanok
  \uses{def:arithmetic-s12-model, def:arithmetic-t12-models}
  Every $L$-structure that is a model of $S^1_2$ is a model of $T^0_2$.
  \source{\cite{Buss86}.}
  \fidelity{Not checked against the printed source.}
  \leanfile{Sagitta/Map/Arranque.lean}
\end{theorem}

\begin{proof}
  \leanok
  Given a sharply bounded $A(x)$ and $a$, apply PIND for a $\Pi^b_1$ formula to $B(z) = \forall w \le
  a\,(w + z \le a \to A(w) \to A(w + z))$. The step from $\lfloor z/2 \rfloor$ to $z$ uses $B(\lfloor
  z/2 \rfloor)$ twice and, if $z$ is odd, once the induction step of $A$.
\end{proof}

\begin{theorem}[Every model of $T^1_2$ is a model of $S^1_2$]
  \label{thm:arithmetic-t12-contains-s12}
  \lean{Sagitta.Map.Acotada.Arranque.t12ContieneS12}
  \leanok
  \uses{def:arithmetic-s12-model, def:arithmetic-t12-models}
  Every $L$-structure that is a model of $T^1_2$ is a model of $S^1_2$.
  \source{\cite{Buss86}.}
  \fidelity{Not checked against the printed source.}
  \leanfile{Sagitta/Map/Arranque2.lean}
\end{theorem}

\begin{proof}
  \leanok
  BASIC is common to both. For PIND, go down through the high-order bits $\lfloor a/2^j \rfloor$ of
  $a$ one at a time by ordinary induction on a $\Sigma^b_1$ formula, from $0$ to $a$; each is the half
  of the next, so the step of PIND is the step of the induction.
\end{proof}

\section{A sequent calculus for $S^1_2$ with PIND as a rule}
\label{sec:arithmetic-sequents}

Mathlib has no proof calculus. The calculus below has the axioms of BASIC as initial sequents and
PIND as an inference rule; its derivations are data.

\begin{definition}[The sequent calculus]
  \label{def:arithmetic-deriva}
  \lean{Sagitta.Map.Acotada.Deriva, Sagitta.Map.Acotada.Vale}
  \leanok
  \uses{def:arithmetic-basic, def:arithmetic-sigma1}
  \own{}
  \infrastructure{}
  A sequent $\Gamma \Rightarrow \Delta$ is a pair of lists of formulas of $L$ with free variables
  indexed by $\mathbb{N}$. A derivation of a sequent is built with the following rules:
  $\varphi \Rightarrow \varphi$; $\Rightarrow \varphi$ for $\varphi$ a quantifier-free instance of an
  axiom of BASIC; $\Rightarrow t = t$; $s = t, A(s) \Rightarrow A(t)$ for $A$ atomic; weakening,
  with contraction and exchange, given by a map from the positions of the premise to those of the
  conclusion; cut; $\bot \Rightarrow$; the rules for $\to$ on both sides; the rules for $\forall$ on
  both sides, the right rule with a variable that does not occur elsewhere; and the rule PIND: from
  $\Gamma, A(\lfloor a/2 \rfloor) \Rightarrow A(a), \Delta$ infer $\Gamma, A(0) \Rightarrow A(t),
  \Delta$, for $A$ a $\Sigma^b_1$ formula, with $a$ not in $\Gamma$, $\Delta$ or the parameters. The
  sequent $\Gamma \Rightarrow \Delta$ \emph{holds} in a structure $M$ under an assignment $v$ when,
  if every formula of $\Gamma$ is true under $v$, some formula of $\Delta$ is.
  \leanfile{Sagitta/Map/Secuentes.lean}
\end{definition}

\begin{theorem}[Soundness of the sequent calculus for the models of $S^1_2$]
  \label{thm:arithmetic-deriva-correcta}
  \lean{Sagitta.Map.Acotada.Deriva.correcta}
  \leanok
  \uses{def:arithmetic-deriva, def:arithmetic-s12-model}
  Let $d$ be a derivation of $\Gamma \Rightarrow \Delta$, $M$ a model of $S^1_2$ and $v$ an
  assignment of the variables to elements of $M$. Then $\Gamma \Rightarrow \Delta$ holds in $M$
  under $v$.
  \source{\cite{Buss98}.}
  \fidelity{Not checked against the printed source.}
  \leanfile{Sagitta/Map/Secuentes.lean}
\end{theorem}

\begin{proof}
  \leanok
  Induction on the derivation. The initial sequents hold because $M$ satisfies BASIC and equality,
  the rule PIND holds because $M$ satisfies PIND for $\Sigma^b_1$ formulas, and the other rules hold
  by the semantics of the connectives.
\end{proof}

\begin{definition}[Free-cut free derivations]
  \label{def:arithmetic-free-cut}
  \lean{Sagitta.Map.Acotada.Deriva.sinCortesLibres, Sagitta.Map.Acotada.Deriva.formulasDeCorte}
  \leanok
  \uses{def:arithmetic-deriva}
  \own{}
  \infrastructure{}
  An occurrence of a formula in a sequent of a derivation is \emph{anchored} when it is a direct
  descendant of a formula of a non-logical initial sequent (BASIC or equality), or of a principal
  formula of an inference PIND. A cut is anchored when its cut formula is atomic and both of its
  occurrences above the cut are anchored, or when it is not atomic and at least one of them is
  anchored. A derivation is free-cut free when every cut in it is anchored. The cut formulas of a
  derivation are the formulas of its cut inferences.
  \leanfile{Sagitta/Map/Secuentes.lean}
\end{definition}

\begin{theorem}[In a free-cut free derivation every cut formula is $\Sigma^b_1$]
  \label{thm:arithmetic-cut-sigma1}
  \lean{Sagitta.Map.Acotada.Deriva.formulasDeCorte_sigma1}
  \leanok
  \uses{def:arithmetic-free-cut, def:arithmetic-sigma1}
  Let $d$ be a derivation of $\Gamma \Rightarrow \Delta$ that is free-cut free. Then every cut
  formula $\varphi$ of $d$ is a $\Sigma^b_1$ formula.
  \source{\cite{Buss98}.}
  \fidelity{Not checked against the printed source.}
  The elimination of free cuts, which would say that every derivable sequent has such a derivation,
  is not proved here; it is stated in the file as the open proposition
  \texttt{EliminacionDeCortesLibres}.
  \leanfile{Sagitta/Map/Secuentes.lean}
\end{theorem}

\begin{proof}
  \leanok
  Induction on $d$. At a cut, the cut formula is atomic, and then it is sharply bounded, or it has an
  anchored occurrence, which descends from an axiom of BASIC or equality or from a principal formula
  of PIND; anchored occurrences are $\Sigma^b_1$.
\end{proof}

\section{The completeness theorem}
\label{sec:arithmetic-completeness}

Provability of a sentence by models is shown to be the same as derivability in the calculus of the
previous section from finitely many axioms of $S^1_2$.

\begin{definition}[Provability by models and derivability from $S^1_2$]
  \label{def:arithmetic-proves}
  \lean{Sagitta.Map.Acotada.Demuestra, Sagitta.Map.Acotada.DerivaS12, Sagitta.Map.Acotada.Der}
  \leanok
  \uses{def:arithmetic-s12-model, def:arithmetic-deriva}
  \own{}
  \infrastructure{}
  $S^1_2$ \emph{proves} a sentence $\sigma$ of $L$ when $\sigma$ is true in every model of $S^1_2$. A
  sequent $\Gamma \Rightarrow \Delta$ is \emph{derivable} when there is a derivation of it. The
  sentence $\sigma$ is derivable from $S^1_2$ when, for some finite list of axioms of the theory
  $S^1_2$, regarded as formulas, the sequent that has these axioms as antecedent and $\sigma$ as
  the only formula of the succedent is derivable.
  \leanfile{Sagitta/Map/Completitud/Teorema.lean}
\end{definition}

\begin{theorem}[Completeness of the sequent calculus, for any theory]
  \label{thm:arithmetic-completeness-theory}
  \lean{Sagitta.Map.Acotada.completitud_teoria}
  \leanok
  \uses{def:arithmetic-deriva}
  Let $\mathrm{Th}$ be a theory of $L$ and $\Gamma$, $\Delta$ lists of formulas with free variables in
  $\mathbb{N}$. Assume that, for every structure $M$ that satisfies $\mathrm{Th}$ and every
  assignment $v : \mathbb{N} \to M$, the sequent $\Gamma \Rightarrow \Delta$ holds in $M$ under
  $v$. Then there is a list $\mathrm{As}$ of formulas, each the formula of a sentence of
  $\mathrm{Th}$, such that $\mathrm{As}, \Gamma \Rightarrow \Delta$ is derivable.
  \source{\cite{Godel30, Henkin49}.}
  \fidelity{Not checked against the printed source.}
  \leanfile{Sagitta/Map/Completitud/Teorema.lean}
\end{theorem}

\begin{proof}
  \leanok
  By contradiction. If no finite list of axioms of $\mathrm{Th}$ derives the sequent, the axioms, the
  antecedent and the negations of the succedent are consistent. Lindenbaum's lemma extends them to a
  maximal Henkin set, and the term model of this set satisfies $\mathrm{Th}$ and falsifies the
  sequent.
\end{proof}

\begin{theorem}[Provability by models is derivability in the sequent calculus]
  \label{thm:arithmetic-demuestra-iff-deriva}
  \lean{Sagitta.Map.Acotada.demuestra_iff_deriva}
  \leanok
  \uses{def:arithmetic-proves, thm:arithmetic-completeness-theory, thm:arithmetic-deriva-correcta}
  For every sentence $\sigma$ of $L$, $S^1_2$ proves $\sigma$ if and only if $\sigma$ is derivable
  from $S^1_2$.
  \source{\cite{Godel30, Henkin49}.}
  \fidelity{Not checked against the printed source.}
  \leanfile{Sagitta/Map/Completitud/Teorema.lean}
\end{theorem}

\begin{proof}
  \leanok
  The forward direction is Theorem~\ref{thm:arithmetic-completeness-theory} applied to $S^1_2$, the
  empty antecedent and the succedent $\sigma$. The converse is soundness
  (Theorem~\ref{thm:arithmetic-deriva-correcta}) applied to a model of $S^1_2$, which satisfies the
  axioms in the antecedent.
\end{proof}

\section{Cobham functions}
\label{sec:arithmetic-cobham}

The programs of Cobham are a syntax for the functions of polynomial time built by bounded recursion on
notation. They live in a model of their own, a type of programs with an evaluation function, and are
tied to machines only by the last entry of this section.

\begin{definition}[Cobham programs and their evaluation]
  \label{def:arithmetic-cob}
  \lean{Sagitta.Map.Cobham.Cob, Sagitta.Map.Cobham.Cob.eval, Sagitta.Map.Cobham.notacion}
  \leanok
  \source{\cite{Cobham65}.}
  \fidelity{Not checked against the printed source.}
  A Cobham program with $n$ arguments is built from: the constant $0$ and the projections $v \mapsto
  v_i$ ($n$ arguments); the successor, the sum, the truncated difference $x \dot- y$, the product,
  $\lfloor x/2 \rfloor$, the length $|x|$ and $x \# y = 2^{|x||y|}$; the conditional, which gives $y$
  if $x = 0$ and $z$ otherwise; the composition $f(g_1(v), \dots, g_m(v))$; and the bounded recursion
  on notation $\mathrm{rn}(g, h, k)$ with $g$ of $n$ arguments, $h$ of $n + 2$ and $k$ of $n + 1$. The
  evaluation is the function of the arguments given by these clauses; for the recursion it is the
  function $F$ with $F(0, \bar y) = g(\bar y)$ and, for $x > 0$,
  \[
    F(x, \bar y) = \min\bigl(h(x, F(\lfloor x/2 \rfloor, \bar y), \bar y),\; k(x, \bar y)\bigr).
  \]
  \leanfile{Sagitta/Map/Cobham.lean}
\end{definition}

\begin{theorem}[Cobham programs run in polynomial time on multi-tape machines]
  \label{thm:arithmetic-cob-polytime}
  \lean{Sagitta.Map.Cobham.cob_polinomial}
  \leanok
  \uses{def:arithmetic-cob, def:polytime}
  For every $n$ and every Cobham program $c$ with $n$ arguments, the function $c.\mathrm{eval} :
  (\mathrm{Fin}\ n \to \mathbb{N}) \to \mathbb{N}$ is computable in polynomial time by a multi-tape
  machine, when a tuple of $n$ numbers is encoded as the tuple of their binary encodings and the
  result as a binary number.
  \source{\cite{Cobham65}.}
  \fidelity{Not checked against the printed source.}
  \leanfile{Sagitta/Map/TiempoCobham.lean}
\end{theorem}

\begin{proof}
  \leanok
  Induction on $c$. The leaves are solved by the tactic \texttt{polytime} of the machine library, the
  composition is a composition of machines with a tuple, and the bounded recursion on notation is
  computed from the top, by a loop whose state grows by at most a polynomial in each step.
\end{proof}

\section{Functions definable in $S^1_2$ and the witnessing lemma}
\label{sec:arithmetic-witness}

A function is definable in $S^1_2$ when a $\Sigma^b_1$ formula defines it and $S^1_2$ proves that the
formula is total. The functions that satisfy this are those of polynomial time.

\begin{definition}[Provably total $\Sigma^b_1$ formulas and $\Sigma^b_1$-definable functions]
  \label{def:arithmetic-definable}
  \lean{Sagitta.Map.Acotada.DemuestraTotal, Sagitta.Map.Acotada.DefinibleSigma1}
  \leanok
  \uses{def:arithmetic-s12-model, def:arithmetic-sigma1}
  \source{\cite{Buss86}.}
  \fidelity{Not checked against the printed source.}
  A formula $\varphi(x, y)$ of $L$ with two free variables is \emph{provably total in $S^1_2$} when
  there is a term $t(x)$ such that in every model $M$ of $S^1_2$: for every $x \in M$ there is $y \in
  M$ with $y \le t(x)$ and $\varphi(x, y)$; and for all $x, y, z \in M$, if $\varphi(x, y)$ and
  $\varphi(x, z)$ then $y = z$. A function $f : \mathbb{N} \to \mathbb{N}$ is \emph{$\Sigma^b_1$-definable
  in $S^1_2$} when there is a $\Sigma^b_1$ formula $\varphi(x, y)$ that is provably total in
  $S^1_2$ and such that $\varphi(n, f(n))$ is true in $\mathbb{N}$ for every $n$.
  \leanfile{Sagitta/Map/Definibles.lean}
\end{definition}

\begin{theorem}[$\Sigma^b_1$-definable functions are computed by Cobham programs]
  \label{thm:arithmetic-definable-cobham}
  \lean{Sagitta.Map.Cobham.definible_es_cobham}
  \leanok
  \uses{def:arithmetic-definable, def:arithmetic-cob}
  Let $f : \mathbb{N} \to \mathbb{N}$ be $\Sigma^b_1$-definable in $S^1_2$. Then there is a Cobham
  program $c$ with one argument such that $c.\mathrm{eval}(n) = f(n)$ for every $n$.
  \source{\cite{Buss86}.}
  \fidelity{Not checked against the printed source.}
  This is the witnessing theorem of Buss, proved by models and not by cut elimination.
  \leanfile{Sagitta/Map/Conservacion.lean}
\end{theorem}

\begin{proof}
  \leanok
  Represent $\varphi$ and ``there is $y \le t(x)$ with $\varphi$'' by fixed Cobham programs. If some
  program gives a witness in $\mathbb{N}$ for every input, its first half is a $y$ with $\varphi(x,
  y)$, and by the uniqueness proved in $S^1_2$ it is $f(x)$. If no program does, then a saturated
  model of $S^1_2$ with an element $d$ that has no witness is built by Herbrand and compactness; there
  $S^1_2$ gives a witness, a contradiction.
\end{proof}

\begin{theorem}[$\Sigma^b_1$-definable functions are of polynomial time]
  \label{thm:arithmetic-definable-polytime}
  \lean{Sagitta.Map.Cobham.definibleEsPolinomial}
  \leanok
  \uses{thm:arithmetic-definable-cobham, thm:arithmetic-cob-polytime, def:arithmetic-definable}
  % uses-lean: Sagitta.Map.EsPolinomial
  Let $f : \mathbb{N} \to \mathbb{N}$ be $\Sigma^b_1$-definable in $S^1_2$. Then $f$ is computable in
  polynomial time by a Turing machine of Mathlib, with the numbers in binary.
  \source{\cite{Buss86}.}
  \fidelity{Not checked against the printed source.}
  \leanfile{Sagitta/Map/TiempoCobham.lean}
\end{theorem}

\begin{proof}
  \leanok
  By Theorem~\ref{thm:arithmetic-definable-cobham} some Cobham program computes $f$; by
  Theorem~\ref{thm:arithmetic-cob-polytime} it runs in polynomial time, and composing with $n
  \mapsto (n)$ and passing to the machines of Mathlib gives the claim.
\end{proof}

\begin{definition}[The witnessing calculus on $\Sigma^b_1$ formulas in negation normal form]
  \label{def:arithmetic-witness-calculus}
  \lean{Sagitta.Map.Extraccion.Termino, Sagitta.Map.Extraccion.Termino.val, Sagitta.Map.Extraccion.Entorno, Sagitta.Map.Extraccion.ext, Sagitta.Map.Extraccion.Chico, Sagitta.Map.Extraccion.Formula, Sagitta.Map.Extraccion.Formula.Vale, Sagitta.Map.Extraccion.Deriv, Sagitta.Map.Extraccion.Arbol, Sagitta.Map.Extraccion.Arbol.peso, Sagitta.Map.Extraccion.Arbol.valor, Sagitta.Map.Extraccion.pesoL, Sagitta.Map.Extraccion.Prog, Sagitta.Map.Extraccion.Prog.correr, Sagitta.Map.Extraccion.Prog.Limitado}
  \leanok
  \own{}
  \infrastructure{}
  Terms are built from variables, $0$, successor, $+$, $\cdot$, $\lfloor t/2 \rfloor$, $|t|$ and
  $t \# u$; an environment gives a natural number to each variable, and $\mathrm{ext}(a, \rho)$ puts
  $a$ as the variable $0$ in front of $\rho$. Formulas are atoms $t \le u$, $t = u$ and their
  negations, conjunction, disjunction, $\exists x \le t\, \varphi$ and $\forall x \le |t|\,
  \varphi$. A derivation is a sequent calculus with identity, arithmetic facts between atoms,
  structural rules, cut, case split on an atom, the rules for the connectives, the two rules for the
  bounded existential quantifier, the rules for the small universal quantifier, and PIND. A witness
  is a tree; its weight is the number of its nodes plus the number of bits of its numbers. A
  program has one instruction for each rule, and \texttt{correr}$(n, \rho, a)$ runs it on a bound
  $n$ for the lengths of the variables, an environment $\rho$ and a list $a$ of witnesses, and
  returns the position of a formula and its witness. A program is bounded when the bounds of its
  recursions are increasing and polynomial. An environment is small for $n$ when every $\rho(i)$
  has at most $n$ bits.
  \leanfile{Sagitta/Map/Testigo.lean}
\end{definition}

\begin{theorem}[Witnessing lemma, for a calculus of $\Sigma^b_1$ formulas in negation normal form]
  \label{thm:arithmetic-witness-lemma}
  \lean{Sagitta.Map.Extraccion.lema_del_testigo}
  \leanok
  \uses{def:arithmetic-witness-calculus}
  Let $t$ be a term, $\varphi$ a formula, and $d$ a derivation of the sequent $\Rightarrow \exists x
  \le t\, \varphi$. Then there is a bounded program $p$ such that: for all $n$ and every environment
  $\rho$ that is small for $n$, the program $p$ run on $n$, $\rho$ and the empty list of witnesses
  returns a witness whose value $a$ satisfies $a \le t(\rho)$ and $\varphi$ is true in the
  environment $\mathrm{ext}(a, \rho)$; and there is a polynomial $q$ such that for all $n$, all
  small $\rho$ and all lists $a$ of witnesses, the weight of the witness that $p$ returns is at most
  $q(n + \text{weight}(a))$.
  \source{\cite{Buss98}.}
  \fidelity{Not checked against the printed source.}
  The calculus is not the one of Section~\ref{sec:arithmetic-sequents}: its formulas are $\Sigma^b_1$
  by construction, so no cut elimination is needed here.
  \leanfile{Sagitta/Map/Basicos.lean}
\end{theorem}

\begin{proof}
  \leanok
  The program is obtained from $d$ by erasing the formulas, keeping one instruction for each rule;
  that it returns a witness is proved by induction on $d$, and its weight bound by a bound on the
  steps of each rule.
\end{proof}

\section{The easy direction of the theorem of Buss}
\label{sec:arithmetic-easy}

The converse to Theorem~\ref{thm:arithmetic-definable-polytime} is stated with two named hypotheses.

\begin{definition}[$\Sigma^b_1$-definability for several arguments, and Cobham programs without recursion]
  \label{def:arithmetic-definable-n}
  \lean{Sagitta.Map.Acotada.DefinibleCobham.DefinibleN, Sagitta.Map.Acotada.DefinibleCobham.SinRn}
  \leanok
  \uses{def:arithmetic-cob, def:arithmetic-definable}
  \own{}
  \infrastructure{}
  A function $f : (\mathrm{Fin}\ n \to \mathbb{N}) \to \mathbb{N}$ is $\Sigma^b_1$-definable in
  $S^1_2$ when there are a $\Sigma^b_1$ formula $\varphi(x_1, \dots, x_n, y)$ and a term $t(\bar x)$
  such that in every model of $S^1_2$ every $\bar x$ has some $y \le t(\bar x)$ with $\varphi(\bar x,
  y)$, this $y$ is unique, and $\varphi(v, f(v))$ is true in $\mathbb{N}$ for every $v$. A Cobham
  program has no recursion when it is built with the basic programs and composition only.
  \leanfile{Sagitta/Map/DefinibleCobham.lean}
\end{definition}

\begin{definition}[Hypothesis: polynomial-time functions are Cobham programs]
  \label{def:arithmetic-hyp-pol-cobham}
  \lean{Sagitta.Map.Acotada.DefinibleCobham.PolinomialEsCobham}
  \leanok
  \uses{def:arithmetic-cob}
  % uses-lean: Sagitta.Map.EsPolinomial
  The statement that for every $f : \mathbb{N} \to \mathbb{N}$ computable in polynomial time by a
  Turing machine of Mathlib, with the numbers in binary, there is a Cobham program $c$ with one
  argument such that $c.\mathrm{eval}(n) = f(n)$ for every $n$. It is a definition of type
  \texttt{Prop}; it is a known theorem of Cobham, and it is not proved in this repository.
  \source{\cite{Cobham65}.}
  \fidelity{Not checked against the printed source.}
  \leanfile{Sagitta/Map/DefinibleCobham.lean}
\end{definition}

\begin{definition}[Hypothesis: bounded recursion on notation preserves definability]
  \label{def:arithmetic-hyp-recursion}
  \lean{Sagitta.Map.Acotada.DefinibleCobham.RecursionDefinible}
  \leanok
  \uses{def:arithmetic-definable-n, def:arithmetic-cob}
  The statement that for all $n$, $g$, $h$, $k$ Cobham programs with $n$, $n + 2$ and $n + 1$
  arguments, if $g.\mathrm{eval}$, $h.\mathrm{eval}$ and $k.\mathrm{eval}$ are $\Sigma^b_1$-definable
  in $S^1_2$, then so is the evaluation of $\mathrm{rn}(g, h, k)$. It is a definition of type
  \texttt{Prop}; it is a known theorem of Buss and it is not proved in this repository.
  \source{\cite{Buss86}.}
  \fidelity{Not checked against the printed source.}
  \leanfile{Sagitta/Map/DefinibleCobham.lean}
\end{definition}

\begin{theorem}[Cobham programs without recursion are $\Sigma^b_1$-definable]
  \label{thm:arithmetic-sinrn-definable}
  \lean{Sagitta.Map.Acotada.DefinibleCobham.sinRn_definible}
  \leanok
  \uses{def:arithmetic-definable-n}
  Let $c$ be a Cobham program with $n$ arguments that has no recursion. Then $c.\mathrm{eval}$ is
  $\Sigma^b_1$-definable in $S^1_2$.
  \source{\cite{Buss86}.}
  \fidelity{Not checked against the printed source.}
  \leanfile{Sagitta/Map/DefinibleCobham.lean}
\end{theorem}

\begin{proof}
  \leanok
  Induction on the proof that $c$ has no recursion. The basic programs are terms, the truncated
  difference or the conditional, each definable by a $\Sigma^b_1$ formula, and the composition of
  definable functions is definable.
\end{proof}

\begin{theorem}[The theorem of Buss, from two hypotheses]
  \label{thm:arithmetic-buss}
  \lean{Sagitta.Map.Acotada.DefinibleCobham.teoremaDeBuss_de}
  \leanok
  \uses{def:arithmetic-hyp-pol-cobham, def:arithmetic-hyp-recursion, thm:arithmetic-definable-polytime, def:arithmetic-definable}
  % uses-lean: Sagitta.Map.EsPolinomial
  Assume that polynomial-time functions are Cobham programs (Definition~\ref{def:arithmetic-hyp-pol-cobham})
  and that bounded recursion on notation preserves definability
  (Definition~\ref{def:arithmetic-hyp-recursion}). Then for every $f : \mathbb{N} \to \mathbb{N}$, $f$
  is $\Sigma^b_1$-definable in $S^1_2$ if and only if $f$ is computable in polynomial time.
  \source{\cite{Buss86}.}
  \fidelity{Not checked against the printed source.}
  \leanfile{Sagitta/Map/DefinibleCobham.lean}
\end{theorem}

\begin{proof}
  \leanok
  The forward direction is Theorem~\ref{thm:arithmetic-definable-polytime}. For the converse, the
  first hypothesis gives a Cobham program for $f$; with the second, every Cobham program is
  definable, by induction on the program, and a definable program with one argument gives a
  $\Sigma^b_1$-definition of the function.
\end{proof}

\section{The theorem of Parikh}
\label{sec:arithmetic-parikh}

Parikh's theorem is stated for $S^1_2$ with $\Sigma^b_1$ formulas and unique outputs, and for $T^1_2$
with bounded formulas and no uniqueness.

\begin{theorem}[The theorem of Parikh for $S^1_2$, for $\Sigma^b_1$ formulas with a unique output]
  \label{thm:arithmetic-parikh}
  \lean{Sagitta.Map.Acotada.Parikh.parikh}
  \leanok
  \uses{def:arithmetic-sigma1, def:arithmetic-definable, def:arithmetic-s12-model}
  Let $\varphi(x, y)$ be a $\Sigma^b_1$ formula with two free variables. Assume that in every model
  $M$ of $S^1_2$ and for every $x \in M$ there is exactly one $y \in M$ with $\varphi(x, y)$. Then
  there is a term $t(x)$ such that in every model $M$ of $S^1_2$: for every $x \in M$ there is $y \le
  t(x)$ with $\varphi(x, y)$, and for all $x, y, z \in M$, if $\varphi(x, y)$ and $\varphi(x, z)$ then
  $y = z$.
  \source{\cite{Parikh71}.}
  \fidelity{Not checked against the printed source.}
  \leanfile{Sagitta/Map/Parikh.lean}
\end{theorem}

\begin{proof}
  \leanok
  If no term bounds the output, compactness gives a model of $S^1_2$ with an element $c$ for which no
  term does. The elements bounded by a term in $c$ form an initial segment that is closed under the
  functions, because they are monotone in every model of $S^1_2$; it is a model of BASIC because BASIC
  is universal, and of PIND because $\Sigma^b_1$ formulas say the same in the segment as in the
  model. In the segment $S^1_2$ gives an output for $c$, bounded by a term, a contradiction.
\end{proof}

\begin{definition}[Bounded formulas]
  \label{def:arithmetic-bounded}
  \lean{Sagitta.Map.Borde.OscilacionParikh.EsAcotada}
  \leanok
  \uses{def:arithmetic-basic}
  \own{}
  \infrastructure{}
  The bounded formulas of $L$ are the smallest class that contains $\bot$ and the atomic formulas, is
  closed under implication, and is closed under $\forall x \le t$ and $\exists x \le t$ for every
  term $t$.
  \leanfile{Sagitta/Map/Borde/OscilacionParikh.lean}
\end{definition}

\begin{theorem}[The theorem of Parikh for $T^1_2$, for bounded formulas without uniqueness]
  \label{thm:arithmetic-parikh-t12}
  \lean{Sagitta.Map.Borde.ParikhT12.parikh_existe_T12}
  \leanok
  \uses{def:arithmetic-bounded, def:arithmetic-t12-models}
  Let $\varphi(x, y)$ be a bounded formula with two free variables. Assume that in every model $M$ of
  $T^1_2$ and for every $x \in M$ there is $y \in M$ with $\varphi(x, y)$. Then there is a term $t(x)$
  such that in every model $M$ of $T^1_2$ and for every $x \in M$ there is $y \in M$ with $y \le t(x)$
  and $\varphi(x, y)$.
  \source{\cite{Parikh71}.}
  \fidelity{Not checked against the printed source.}
  \leanfile{Sagitta/Map/Borde/ParikhT12.lean}
\end{theorem}

\begin{proof}
  \leanok
  It is the argument of the previous theorem: compactness, the initial segment of the elements
  bounded by a term in a counterexample element, which is a model of $T^1_2$ because the formulas
  are bounded and the functions are monotone (Theorem~\ref{thm:arithmetic-t12-contains-s12}), and
  the output that the theory gives in the segment.
\end{proof}

\section{Decision by two bounded formulas, and $\cP = \cNP$}
\label{sec:arithmetic-delta}

A problem is a function from lists of bits to booleans. A string $w$ is read as a number, and a
formula with one free variable defines a problem when it is true of the number of $w$ exactly for
the strings of the problem.

\begin{definition}[A formula defines a problem]
  \label{def:arithmetic-defines}
  \lean{Sagitta.Map.Acotada.DefinidoPor, Sagitta.Map.Acotada.numeroDe}
  \leanok
  \own{}
  \infrastructure{}
  % uses-lean: Sagitta.Map.Problema
  For a list of bits $w$, the number of $w$ is the number whose binary expansion is the digit $1$
  followed by the digits of $w$, the first bit of $w$ being the least significant. A formula
  $\varphi(x)$ of $L$ with one free variable \emph{defines} the problem $P$ when, for every list of
  bits $w$, $P(w) = \mathrm{true}$ if and only if $\varphi$ is true in $\mathbb{N}$ of the number of
  $w$.
  \leanfile{Sagitta/Map/SigmaNP.lean}
\end{definition}

\begin{definition}[Hypothesis: $S^1_2$ decides between two formulas]
  \label{def:arithmetic-hyp-decides}
  \lean{Sagitta.Map.DeltaUno.DemuestraDisyuncion}
  \leanok
  \uses{def:arithmetic-s12-model}
  For formulas $\varphi(x)$ and $\psi(x)$ with one free variable, the statement that in every model
  $M$ of $S^1_2$ and for every $x \in M$, $\varphi(x)$ or $\psi(x)$ holds. It is a definition of
  type \texttt{Prop}.
  \source{\cite{Buss86}.}
  \fidelity{Not checked against the printed source.}
  \leanfile{Sagitta/Map/DeltaUno.lean}
\end{definition}

\begin{theorem}[What $S^1_2$ decides with two $\Sigma^b_1$ formulas is in $\cP$]
  \label{thm:arithmetic-delta-p}
  \lean{Sagitta.Map.DeltaUno.deltaUno_en_P}
  \leanok
  \uses{def:arithmetic-hyp-decides, def:arithmetic-defines, def:arithmetic-sigma1, thm:arithmetic-definable-cobham, thm:arithmetic-cob-polytime}
  % uses-lean: Sagitta.Map.P, Sagitta.Map.Problema
  Let $\varphi(x)$ and $\psi(x)$ be $\Sigma^b_1$ formulas with one free variable. Assume that
  $S^1_2$ decides between $\varphi$ and $\psi$ (Definition~\ref{def:arithmetic-hyp-decides}). Let
  $\mathit{Pr}$ be a problem such that $\varphi$ defines $\mathit{Pr}$ and $\psi$ defines the
  problem $w \mapsto \neg \mathit{Pr}(w)$. Then $\mathit{Pr}$ is decided by a Turing machine of
  Mathlib in polynomial time.
  \source{\cite{Buss86}.}
  \fidelity{Not checked against the printed source.}
  \leanfile{Sagitta/Map/DeltaUno.lean}
\end{theorem}

\begin{proof}
  \leanok
  By the witnessing argument there is a Cobham program $c$ such that for every $n$, either
  $c(n) = 0$ and $\varphi(n)$, or $c(n) = 1$ and $\psi(n)$. Then $\mathit{Pr}(w)$ is true if and only
  if $c$ of the number of $w$ is $0$, which is polynomial time by the machines for Cobham programs.
\end{proof}

\begin{definition}[Hypothesis: $S^1_2$ decides SAT]
  \label{def:arithmetic-hyp-decides-sat}
  \lean{Sagitta.Map.CookS12.S12DecideSat}
  \leanok
  \uses{def:arithmetic-hyp-decides, def:arithmetic-defines, def:arithmetic-sigma1}
  The statement that there are $\Sigma^b_1$ formulas $\varphi(x)$ and $\psi(x)$ with one free
  variable such that $\varphi$ defines the problem SAT, $\psi$ defines its complement, and $S^1_2$
  decides between $\varphi$ and $\psi$. Here SAT is the problem of the strings that encode a
  satisfiable CNF formula. It is a definition of type \texttt{Prop}.
  \own{}
  \leanfile{Sagitta/Map/CookS12.lean}
\end{definition}

\begin{theorem}[If $S^1_2$ decides SAT, then $\cP = \cNP$]
  \label{thm:arithmetic-decides-sat-pnp}
  \lean{Sagitta.Map.CookS12.pIgualNP_si_s12_decide_sat}
  \leanok
  \uses{def:arithmetic-hyp-decides-sat, thm:arithmetic-delta-p}
  % uses-lean: Sagitta.Map.P, Sagitta.Map.NP
  Assume that $S^1_2$ decides SAT (Definition~\ref{def:arithmetic-hyp-decides-sat}). Then $\cP =
  \cNP$, as sets of problems.
  \source{\cite{Buss86}.}
  \fidelity{Not checked against the printed source.}
  \leanfile{Sagitta/Map/CookS12.lean}
\end{theorem}

\begin{proof}
  \leanok
  By Theorem~\ref{thm:arithmetic-delta-p}, SAT is in $\cP$. Since SAT is $\cNP$-hard (the theorem of
  Cook and Levin of the machine library), $\cNP \subseteq \cP$, and $\cP \subseteq \cNP$ is
  known.
\end{proof}

\section{Bounded formulas and the classes $\cNP$ and $\cconP$}
\label{sec:arithmetic-sigma-np}

The classes of problems that $\Sigma^b_1$ and $\Pi^b_1$ formulas define are those of $\cNP$ and
$\cconP$.

\begin{theorem}[The problems defined by $\Sigma^b_1$ formulas are those of $\cNP$]
  \label{thm:arithmetic-sigma-np}
  \lean{Sagitta.Map.SigmaFagin.sigmaIgualNP, Sagitta.Map.SigmaFagin.sigmaSubNP, Sagitta.Map.SigmaFagin.npSubSigma}
  \leanok
  \uses{def:arithmetic-defines, def:arithmetic-sigma1, thm:arithmetic-cob-polytime}
  % uses-lean: Sagitta.Map.NP, Sagitta.Map.Problema
  For every problem $P$, $P \in \cNP$ if and only if there is a $\Sigma^b_1$ formula $\varphi(x)$ with
  one free variable that defines $P$. The direction from the formula to $\cNP$ is
  \texttt{sigmaSubNP}, and the direction from $\cNP$ to the formula is \texttt{npSubSigma}.
  \own{}
  \leanfile{Sagitta/Map/SigmaFagin/Palabras.lean}
\end{theorem}

\begin{proof}
  \leanok
  For $\Sigma^b_1 \subseteq \cNP$, a $\Sigma^b_1$ formula is represented in $\mathbb{N}$ as ``there is
  $y \le t$ with $\delta = 0$'' for Cobham programs $\delta$ and $t$, and the verifier computes
  $\delta$, in polynomial time by Theorem~\ref{thm:arithmetic-cob-polytime}. For $\cNP \subseteq
  \Sigma^b_1$, the language of the query of $P$ is in $\cNP$ in the machine library, whose Fagin
  theorem for ordered structures gives an existential second-order sentence, which is translated into
  a $\Sigma^b_1$ formula in which the input is the free variable.
\end{proof}

\begin{theorem}[$\cNP \ne \cconP$ as a question about formulas]
  \label{thm:arithmetic-np-conp-formulas}
  \lean{Sagitta.Map.SigmaFagin.npDistintoCoNP_iff_formulas}
  \leanok
  \uses{thm:arithmetic-sigma-np, def:arithmetic-defines, def:arithmetic-sigma1}
  % uses-lean: Sagitta.Map.NPDistintoCoNP, Sagitta.Map.Problema
  $\cNP \ne \cconP$ if and only if there is a problem $P$ that is defined by a $\Pi^b_1$ formula with
  one free variable and that is defined by no $\Sigma^b_1$ formula with one free variable.
  \own{}
  \leanfile{Sagitta/Map/SigmaFagin/Palabras.lean}
\end{theorem}

\begin{proof}
  \leanok
  It is \texttt{npDistintoCoNP\_iff} with Theorem~\ref{thm:arithmetic-sigma-np}: the $\Sigma^b_1$
  formulas give $\cNP$, their duals the $\Pi^b_1$ formulas give $\cconP$, and the two classes are
  different if and only if some problem is in one and not in the other.
\end{proof}

\section{Student and teacher witnessing for a universal theory of Cobham programs}
\label{sec:arithmetic-kpt}

The result is stated for the universal theory of the true universal sentences of the language that has
one function symbol for each Cobham program, and not for $S^1_2$ or $T^1_2$.

\begin{definition}[The theory of Cobham programs]
  \label{def:arithmetic-tc}
  \lean{Sagitta.Map.Cobham.LC, Sagitta.Map.Cobham.TC}
  \leanok
  \uses{def:arithmetic-cob}
  \own{}
  \infrastructure{}
  The language $L_C$ has a function symbol for each Cobham program, with the arity of the program,
  and no relation symbol. The structure $\mathbb{N}$ interprets each symbol as the evaluation of the
  program. The theory $TC$ is the set of universal sentences of $L_C$ that are true in $\mathbb{N}$.
  \leanfile{Sagitta/Map/TeoriaCobham.lean}
\end{definition}

\begin{definition}[Hypothesis: $TC$ proves $\forall x\, \exists y\, \forall z\, \varphi$]
  \label{def:arithmetic-hyp-aea}
  \lean{Sagitta.Map.Cobham.DemuestraAEA}
  \leanok
  \uses{def:arithmetic-tc}
  For a Cobham program $\chi$ with three arguments, the statement that in every $L_C$-structure $M$
  that satisfies $TC$ and for every $x \in M$ there is $y \in M$ such that for every $z \in M$ the
  value of the symbol $\chi$ on $(x, y, z)$ is the value of the constant program $0$. The formula
  $\varphi$ is ``$\chi$ gives $0$''. It is a definition of type \texttt{Prop}.
  \source{\cite{KPT91}.}
  \fidelity{Not checked against the printed source.}
  \leanfile{Sagitta/Map/KPT.lean}
\end{definition}

\begin{definition}[The student wins in $k$ rounds]
  \label{def:arithmetic-student}
  \lean{Sagitta.Map.Cobham.GanaAlumno, Sagitta.Map.Cobham.previas}
  \leanok
  \uses{def:arithmetic-cob}
  \own{}
  \infrastructure{}
  For a list of counterexamples $z_0, \dots, z_{k-1}$ and a round $i$, the list $\mathrm{previas}(z,
  i)$ has $z_j$ at the places $j < i$ and $0$ elsewhere. For a Cobham program $\chi$ with three
  arguments and programs $f_0, \dots, f_{k-1}$ with $k + 1$ arguments, the student wins in $k$ rounds
  when for every $x \in \mathbb{N}$ and every list $z_0, \dots, z_{k-1}$ of natural numbers there is
  a round $i < k$ with $\chi(x,\ f_i(x, \mathrm{previas}(z, i)),\ z_i) = 0$.
  \leanfile{Sagitta/Map/KPT.lean}
\end{definition}

\begin{theorem}[Student-teacher witnessing for the universal theory $TC$]
  \label{thm:arithmetic-kpt}
  \lean{Sagitta.Map.Cobham.kpt, Sagitta.Map.Cobham.kpt_iff}
  \leanok
  \uses{def:arithmetic-hyp-aea, def:arithmetic-student}
  Let $\chi$ be a Cobham program with three arguments. If $TC$ proves $\forall x\, \exists y\,
  \forall z\, (\chi(x, y, z) = 0)$ (Definition~\ref{def:arithmetic-hyp-aea}), then there are $k$ and
  Cobham programs $f_0, \dots, f_{k-1}$ with which the student wins in $k$ rounds. Conversely, if
  there are such $k$ and programs, $TC$ proves it.
  \source{\cite{KPT91}.}
  \fidelity{Not checked against the printed source.}
  The statement is for the universal theory $TC$, not for $S^1_2$ or $T^1_2$.
  \leanfile{Sagitta/Map/KPT.lean}
\end{theorem}

\begin{proof}
  \leanok
  Enumerate all programs, each infinitely often; round $j$ of the teacher says that the $j$-th
  program of the student fails. If the student never wins in $\mathbb{N}$, every finite set of rounds
  has a model in $\mathbb{N}$, so by compactness all of them have a model. Its closure is again a model
  of the universal theory $TC$, in which every element is a program of the constants, and the
  hypothesis gives a $y$ that is one of the programs, which some round refutes. The converse
  direction is a direct proof from the disjunction of the rounds.
\end{proof}

\section{Statements that $S^1_2$ and $T^1_2$ do not prove if circuit sizes oscillate}
\label{sec:arithmetic-oscillation}

The results of this section are conditional: they assume that a sequence of circuit sizes oscillates
between polynomial and superpolynomial. This is not proved, and not refuted, here.

\begin{definition}[Hypothesis: a sequence oscillates]
  \label{def:arithmetic-hyp-oscillation}
  \lean{Sagitta.Map.Borde.Tricotomia.Bajo, Sagitta.Map.Borde.Tricotomia.Sobre, Sagitta.Map.Borde.Tricotomia.Oscila}
  \leanok
  \own{}
  A function $f : \mathbb{N} \to \mathbb{N}$ is \emph{bounded by a polynomial} when there are $c$ and
  $k$ with $f(n) \le c\,(n + 1)^k$ for every $n$. It is \emph{above every polynomial} when for all $c$
  and $k$, $c\,(n + 1)^k < f(n)$ for all large enough $n$. It \emph{oscillates} when it is neither
  bounded by a polynomial nor above every polynomial. The statement that a given function oscillates
  is a definition of type \texttt{Prop}.
  \leanfile{Sagitta/Map/Borde/Tricotomia.lean}
\end{definition}

\begin{definition}[Good lengths of a problem]
  \label{def:arithmetic-good-lengths}
  \lean{Sagitta.Map.Borde.OscilacionParikh.Buenos}
  \leanok
  \own{}
  % uses-lean: Sagitta.Map.Tiempo.Finitud.sucNat, Sagitta.Map.Tiempo.Relleno.satMarco
  Let $s_P(m)$ be the least number of gates of a Boolean circuit that decides the problem $P$ on the
  strings of length $m$. For a problem $P$ and numbers $c$ and $k$, the relation $\mathrm{Buenos}(P,
  c, k)$ between $x$ and $y$ says that $|x| \le |y|$ and $s_P(|y|) \le c\,(|y| + 1)^k$. The problem
  $\mathrm{satMarco}$ of the next theorems is SAT on framed inputs: the input is unframed, which
  reads it up to the first marker, what is read is asked to SAT, and everything after the frame is
  ignored.
  \leanfile{Sagitta/Map/Borde/OscilacionParikh.lean}
\end{definition}

\begin{definition}[Hypothesis: a sound theory with the Parikh property]
  \label{def:arithmetic-hyp-parikh-sound}
  \lean{Sagitta.Map.Borde.OscilacionParikh.Enunciado, Sagitta.Map.Borde.OscilacionParikh.Verdadero, Sagitta.Map.Borde.OscilacionParikh.PolLargo, Sagitta.Map.Borde.OscilacionParikh.CotaPol, Sagitta.Map.Borde.OscilacionParikh.Acotar, Sagitta.Map.Borde.OscilacionParikh.ParikhSana}
  \leanok
  \own{}
  A statement is a relation $\psi(x, y)$ between natural numbers; it stands for $\forall x\, \exists y\,
  \psi(x, y)$, and it is true when this holds in $\mathbb{N}$. A function $\beta$ has polynomial length
  when $|\beta(x)| \le c\,(|x| + 1)^d$ for some $c$, $d$ and all $x$. The statement $\psi$ with the
  output bounded by $\beta$ is the relation $y \le \beta(x) \wedge \psi(x, y)$. For sets $C$ and
  $\mathrm{Th}$ of statements, $(C, \mathrm{Th})$ is a \emph{sound theory with the Parikh property}
  when every statement of $\mathrm{Th}$ is true, and for every $\psi \in \mathrm{Th}$ with $\psi \in
  C$ there is a function $\beta$ of polynomial length such that $\psi$ with the output bounded by
  $\beta$ is in $\mathrm{Th}$. It is a definition of type \texttt{Prop}.
  \leanfile{Sagitta/Map/Borde/OscilacionParikh.lean}
\end{definition}

\begin{theorem}[If framed SAT oscillates in circuits, no sound theory with the Parikh property proves that its good lengths do not end]
  \label{thm:arithmetic-oscillation-unprovable}
  \lean{Sagitta.Map.Borde.OscilacionParikh.oscila_no_se_demuestra}
  \leanok
  \uses{def:arithmetic-hyp-parikh-sound, def:arithmetic-hyp-oscillation, def:arithmetic-good-lengths}
  % uses-lean: Sagitta.Map.Tiempo.Finitud.sucNat, Sagitta.Map.Tiempo.Relleno.satMarco
  Let $C$ and $\mathrm{Th}$ be sets of statements such that $(C, \mathrm{Th})$ is a sound theory
  with the Parikh property. Assume that the sequence $m \mapsto s_{\mathrm{satMarco}}(m)$ of least
  circuit sizes oscillates. Let $c$ and $k$ be such that $\mathrm{Buenos}(\mathrm{satMarco}, c,
  k) \in C$. Then $\mathrm{Buenos}(\mathrm{satMarco}, c, k) \notin \mathrm{Th}$.
  \source{\cite{BKO20}.}
  \fidelity{Not checked against the printed source.}
  \leanfile{Sagitta/Map/Borde/OscilacionParikh.lean}
\end{theorem}

\begin{proof}
  \leanok
  If $\mathrm{Th}$ proved it, the Parikh property gives a bound of polynomial length for the witnesses,
  which is true by soundness; the good lengths are then polynomially dense, and by padding
  $\mathrm{satMarco}$ has circuits of polynomial size, which contradicts that the sequence is not
  bounded by a polynomial.
\end{proof}

\begin{definition}[Universal closure and provability of $\forall x\, \exists y\, \varphi$]
  \label{def:arithmetic-forall-exists}
  \lean{Sagitta.Map.Acotada.cierre, Sagitta.Map.Borde.OscilacionParikh.existeY, Sagitta.Map.Borde.ParikhT12.DemuestraExisteT12}
  \leanok
  \uses{def:arithmetic-t12-models}
  \own{}
  \infrastructure{}
  For a formula $\chi$ with free variables $x_1, \dots, x_n$, its closure is the sentence $\forall
  x_1 \cdots \forall x_n\, \chi$. For a formula $\varphi(x, y)$, $\exists y\, \varphi$ is the formula
  with the free variable $x$. $T^1_2$ proves $\forall x\, \exists y\, \varphi$ when in every model
  $M$ of $T^1_2$ and for every $x \in M$ there is $y \in M$ with $\varphi(x, y)$.
  \leanfile{Sagitta/Map/Borde/OscilacionParikh.lean}
\end{definition}

\begin{theorem}[If framed SAT oscillates in circuits, a true $\forall x\, \exists y\, \varphi$ with $\varphi$ bounded is not derivable in $S^1_2$]
  \label{thm:arithmetic-s12-oscillation}
  \lean{Sagitta.Map.Borde.FormulaBuenos.s12_oscila_sin_hipotesis}
  \leanok
  \uses{def:arithmetic-hyp-oscillation, def:arithmetic-bounded, def:arithmetic-forall-exists, def:arithmetic-proves, thm:arithmetic-oscillation-unprovable}
  % uses-lean: Sagitta.Map.Tiempo.Finitud.sucNat, Sagitta.Map.Tiempo.Relleno.satMarco
  Assume that the sequence $m \mapsto s_{\mathrm{satMarco}}(m)$ of least circuit sizes oscillates.
  Then there is a bounded formula $\varphi(x, y)$ with two free variables such that, for every $x \in
  \mathbb{N}$, there is $y \in \mathbb{N}$ with $\varphi(x, y)$, and the sentence $\forall x\,
  \exists y\, \varphi(x, y)$ is not derivable from $S^1_2$.
  \source{\cite{BKO20}.}
  \fidelity{Not checked against the printed source.}
  \leanfile{Sagitta/Map/Borde/FormulaBuenosEnsamble.lean}
\end{theorem}

\begin{proof}
  \leanok
  Oscillation gives $c$ and $k$ for which the good lengths do not end. A bounded formula that defines
  $\mathrm{Buenos}(\mathrm{satMarco}, c, k)$ in $\mathbb{N}$ is built from three $\Sigma^b_1$ formulas
  (the shape of a circuit code, its evaluation, and framed SAT) that come from $\Sigma^b_1 = \cNP$.
  The theory $S^1_2$ with its Parikh theorem for bounded formulas is a sound theory with the Parikh
  property, so Theorem~\ref{thm:arithmetic-oscillation-unprovable} applies.
\end{proof}

\begin{theorem}[If framed SAT oscillates in circuits, a true $\forall x\, \exists y\, \varphi$ with $\varphi$ bounded is not proved by $T^1_2$]
  \label{thm:arithmetic-t12-oscillation}
  \lean{Sagitta.Map.Borde.FormulaBuenos.t12_oscila_sin_hipotesis}
  \leanok
  \uses{def:arithmetic-hyp-oscillation, def:arithmetic-bounded, def:arithmetic-forall-exists, thm:arithmetic-parikh-t12, thm:arithmetic-oscillation-unprovable}
  % uses-lean: Sagitta.Map.Tiempo.Finitud.sucNat, Sagitta.Map.Tiempo.Relleno.satMarco
  Assume that the sequence $m \mapsto s_{\mathrm{satMarco}}(m)$ of least circuit sizes oscillates.
  Then there is a bounded formula $\varphi(x, y)$ with two free variables such that, for every $x \in
  \mathbb{N}$, there is $y \in \mathbb{N}$ with $\varphi(x, y)$, and $T^1_2$ does not prove $\forall x\,
  \exists y\, \varphi(x, y)$ (Definition~\ref{def:arithmetic-forall-exists}).
  \source{\cite{BKO20}.}
  \fidelity{Not checked against the printed source.}
  \leanfile{Sagitta/Map/Borde/FormulaBuenosEnsamble.lean}
\end{theorem}

\begin{proof}
  \leanok
  As for $S^1_2$: oscillation gives $c$ and $k$, the formula that defines the good lengths is the same,
  and $T^1_2$ with the Parikh theorem for $T^1_2$ is a sound theory with the Parikh property for the
  bounded formulas.
\end{proof}

\section{Factorization and $S^1_2$}
\label{sec:arithmetic-factor}

The first result is conditional on a hypothesis about the complexity of finding the least prime
factor. The second is a bound on certificates that a verifier can check.

\begin{theorem}[$S^1_2$ does not derive the existence of the least prime factor, if finding it is not polynomial]
  \label{thm:arithmetic-no-factor}
  \lean{Sagitta.Map.Acotada.FactorizacionS12.s12_no_demuestra_factor}
  \leanok
  \uses{def:arithmetic-sigma1, def:arithmetic-proves, def:arithmetic-forall-exists}
  % uses-lean: Sagitta.Map.EsPolinomial, Sagitta.Map.Acotada.existeUnico
  Assume that the function $n \mapsto \mathrm{minFac}(n)$, the least prime factor of $n$, is not
  computable in polynomial time by a Turing machine of Mathlib. Let $\varphi(x, y)$ be a $\Sigma^b_1$
  formula with two free variables such that $\varphi(n, \mathrm{minFac}(n))$ is true in $\mathbb{N}$
  for every $n$. Then the closure of $\exists! y\, \varphi(x, y)$ is not derivable from $S^1_2$.
  \own{}
  \leanfile{Sagitta/Map/Borde/FactorizacionS12.lean}
\end{theorem}

\begin{proof}
  \leanok
  If the sentence were derivable from $S^1_2$, the witnessing theorem of Buss for derivations would
  make $\mathrm{minFac}$ computable in polynomial time, against the assumption.
\end{proof}

\begin{definition}[Pratt certificates and the check of the graph of the least prime factor]
  \label{def:arithmetic-pratt}
  \lean{Sagitta.Map.Acotada.Pratt.Entrada, Sagitta.Map.Acotada.Pratt.entradaOK, Sagitta.Map.Acotada.Pratt.prattOK, Sagitta.Map.Acotada.Pratt.revisaGrafo, Sagitta.Map.Acotada.Pratt.encCertificado}
  \leanok
  \own{}
  \infrastructure{}
  An entry is a triple $(p, a, [q_1, \dots, q_k])$ of a number, a number and a list of numbers. An
  entry is correct within a list $C$ of entries when $p \ge 2$, the product of the $q_i$ is $p - 1$,
  $a^{p-1} \equiv 1 \pmod p$, and every $q_i$ is the first component of some entry of $C$ and
  satisfies $a^{(p-1)/q_i} \not\equiv 1 \pmod p$. A list of entries is a Pratt certificate when every
  entry is correct within it. For numbers $n$ and $m$ and a pair $c = (L, C)$ of a list $L$ of
  numbers and a list $C$ of entries, the check of the graph of the least prime factor is: if $n = 0$,
  that $m = 2$; if $n = 1$, that $m = 1$; otherwise, that the product of $L$ is $n$, $m \in L$, every
  $q \in L$ satisfies $m \le q$ and is the first component of an entry of $C$, and $C$ is a Pratt
  certificate. The encoding of $c$ is a list of bits.
  \leanfile{Sagitta/Map/Borde/Pratt.lean}
\end{definition}

\begin{theorem}[Every number has a short certificate that $m$ is its least prime factor]
  \label{thm:arithmetic-pratt-cert}
  \lean{Sagitta.Map.Acotada.Pratt.certificado_grafo}
  \leanok
  \uses{def:arithmetic-pratt}
  % uses-lean: Sagitta.Map.Acotada.Pratt.encCertificado
  For every natural number $n$ there is a pair $c = (L, C)$, with $L$ a list of numbers and $C$ a list
  of entries, such that the check of the graph of the least prime factor for $n$, $\mathrm{minFac}(n)$
  and $c$ succeeds, and the length of the binary encoding of $c$ is at most $63\,(|n| + 1)^2$, where
  $|n|$ is the number of binary digits of $n$.
  \own{}
  \leanfile{Sagitta/Map/Borde/Pratt.lean}
\end{theorem}

\begin{proof}
  \leanok
  For $n < 2$ the empty pair works. Otherwise let $L$ be the list of prime factors of $n$ with
  repetition; every prime has a Pratt certificate with a bounded number of entries, from the converse
  of the primality criterion of Lucas, and the lengths of the entries are bounded using that $|a| +
  |b| \le |ab| + 1$.
\end{proof}

\section{The second theorem of G\"odel and the theorem of L\"ob, for an abstract system}
\label{sec:arithmetic-godel}

Both theorems are stated for an abstract system of statements with a provability predicate, with the
conditions of derivability as hypotheses. No arithmetic is involved.

\begin{definition}[An abstract system of proofs with an internal provability predicate]
  \label{def:arithmetic-hbl}
  \lean{Sagitta.Map.Tiempo.DiagonalGodel.SistemaHBL, Sagitta.Map.Tiempo.DiagonalGodel.SistemaHBL.caja, Sagitta.Map.Tiempo.DiagonalGodel.SistemaHBL.neg, Sagitta.Map.Tiempo.DiagonalGodel.SistemaHBL.Con}
  \leanok
  \own{}
  \infrastructure{}
  A system $T$ over a type $S$ of statements and a type $C$ of codes has an implication $a \to b$ and
  a statement $\bot$ of $S$; a predicate $\mathrm{Pr}$ on $S$ (provability) that is closed under
  modus ponens and contains the axioms $a \to (b \to a)$ and $(a \to (b \to c)) \to ((a \to b) \to (a
  \to c))$; a code $\ulcorner a \urcorner$ in $C$ for each statement $a$; and a map $\mathrm{Dem}$ from
  $C$ to $S$. Write $\Box a$ for $\mathrm{Dem}(\ulcorner a \urcorner)$, $\neg a$ for $a \to \bot$, and
  $\mathrm{Con}$ for $\neg \Box \bot$.
  \leanfile{Sagitta/Map/DiagonalGodel.lean}
\end{definition}

\begin{definition}[Hypothesis: the derivability conditions]
  \label{def:arithmetic-derivability}
  \lean{Sagitta.Map.Tiempo.DiagonalGodel.SistemaHBL.D1, Sagitta.Map.Tiempo.DiagonalGodel.SistemaHBL.D2, Sagitta.Map.Tiempo.DiagonalGodel.SistemaHBL.D3}
  \leanok
  \uses{def:arithmetic-hbl}
  For a system $T$ as in Definition~\ref{def:arithmetic-hbl}, the three statements: D1, for every $a$,
  if $\mathrm{Pr}(a)$ then $\mathrm{Pr}(\Box a)$; D2, for all $a$ and $b$, $\mathrm{Pr}(\Box(a \to b)
  \to (\Box a \to \Box b))$; D3, for every $a$, $\mathrm{Pr}(\Box a \to \Box\Box a)$. Each is a
  definition of type \texttt{Prop}.
  \source{\cite{Godel31}.}
  \fidelity{Not checked against the printed source.}
  \leanfile{Sagitta/Map/DiagonalGodel.lean}
\end{definition}

\begin{definition}[Hypothesis: a fixed point of the diagonal lemma]
  \label{def:arithmetic-fixed-point}
  \lean{Sagitta.Map.Tiempo.DiagonalGodel.SistemaHBL.PuntoFijoLob, Sagitta.Map.Tiempo.DiagonalGodel.SistemaHBL.PuntoFijo}
  \leanok
  \uses{def:arithmetic-hbl}
  For statements $\varphi$ and $\psi$, $\psi$ is a fixed point for L\"ob of $\varphi$ when $\mathrm{Pr}(\psi
  \to (\Box\psi \to \varphi))$ and $\mathrm{Pr}((\Box\psi \to \varphi) \to \psi)$. A statement $G$ is a
  fixed point of G\"odel when it is a fixed point for L\"ob of $\bot$. Each is a definition of type
  \texttt{Prop}.
  \source{\cite{Godel31}.}
  \fidelity{Not checked against the printed source.}
  \leanfile{Sagitta/Map/DiagonalGodel.lean}
\end{definition}

\begin{theorem}[L\"ob's theorem for an abstract system]
  \label{thm:arithmetic-lob}
  \lean{Sagitta.Map.Tiempo.DiagonalGodel.SistemaHBL.lob}
  \leanok
  \uses{def:arithmetic-derivability, def:arithmetic-fixed-point}
  Let $T$ be a system as in Definition~\ref{def:arithmetic-hbl}. Assume D1, D2 and D3
  (Definition~\ref{def:arithmetic-derivability}). Let $\varphi$ and $\psi$ be statements such that $\psi$
  is a fixed point for L\"ob of $\varphi$ (Definition~\ref{def:arithmetic-fixed-point}). If
  $\mathrm{Pr}(\Box\varphi \to \varphi)$, then $\mathrm{Pr}(\varphi)$.
  \source{\cite{Lob55}.}
  \fidelity{Not checked against the printed source.}
  \leanfile{Sagitta/Map/DiagonalGodel.lean}
\end{theorem}

\begin{proof}
  \leanok
  From D1, D2 and D3 and the first half of the fixed point, $\mathrm{Pr}(\Box\psi \to \Box(\Box\psi
  \to \varphi))$ follows, and with the assumption $\mathrm{Pr}(\Box\psi \to \varphi)$. The second
  half of the fixed point gives $\mathrm{Pr}(\psi)$, D1 gives $\mathrm{Pr}(\Box\psi)$, and modus ponens
  gives $\mathrm{Pr}(\varphi)$.
\end{proof}

\begin{theorem}[The second theorem of G\"odel for an abstract system]
  \label{thm:arithmetic-godel2}
  \lean{Sagitta.Map.Tiempo.DiagonalGodel.SistemaHBL.segundo_teorema}
  \leanok
  \uses{thm:arithmetic-lob, def:arithmetic-derivability, def:arithmetic-fixed-point}
  Let $T$ be a system as in Definition~\ref{def:arithmetic-hbl}. Assume D1, D2 and D3
  (Definition~\ref{def:arithmetic-derivability}). Let $G$ be a fixed point of G\"odel
  (Definition~\ref{def:arithmetic-fixed-point}). If $\mathrm{Pr}(\mathrm{Con})$, then
  $\mathrm{Pr}(\bot)$.
  \source{\cite{Godel31}.}
  \fidelity{Not checked against the printed source.}
  \leanfile{Sagitta/Map/DiagonalGodel.lean}
\end{theorem}

\begin{proof}
  \leanok
  It is Theorem~\ref{thm:arithmetic-lob} with $\varphi = \bot$, since $\mathrm{Con}$ is the statement
  $\Box\bot \to \bot$.
\end{proof}

\section{A propositional mirror of $S^1_2$ and the tautology of circuits}
\label{sec:arithmetic-mirror}

Three results about the translation of bounded formulas into propositional formulas, each in a model
of its own: the formulas of the witnessing calculus, the system of Frege of the machine library, and
Boolean circuits.

\begin{definition}[Translation of sharply bounded formulas into propositional formulas]
  \label{def:arithmetic-translation}
  \lean{Sagitta.Map.Extraccion.EspejoEF.FP, Sagitta.Map.Extraccion.EspejoEF.Bits, Sagitta.Map.Extraccion.EspejoEF.evalB, Sagitta.Map.Extraccion.EspejoEF.EsSb0, Sagitta.Map.Extraccion.EspejoEF.trad}
  \leanok
  \uses{def:arithmetic-witness-calculus}
  \own{}
  \infrastructure{}
  A propositional formula has variables indexed by $\mathbb{N}$ and the connectives $\bot$, $\wedge$,
  $\vee$ and $\to$; $\mathrm{evalB}(\sigma, p)$ is its value under the assignment $\sigma$. A vector
  of bits is a map from $\mathbb{N}$ to propositional formulas, the formula of bit $j$. A formula of
  the witnessing calculus is in the class $\Sigma^b_0$ restricted when its atoms are between simple
  terms (variables, $0$, successor, sum and $\lfloor t/2 \rfloor$), its quantifiers are $\forall x
  \le |t|$ and $\exists x \le |t|$ with $t$ simple, and its connectives are conjunction and
  disjunction. The translation $\mathrm{trad}(\varphi, w, E)$ of a formula, for a width $w$ and a
  vector of bits $E(i)$ for each variable $i$, translates the atoms with adders and comparators on
  the bits, a quantifier $\forall x \le |t|$ as the conjunction over $k = 0, \dots, W$ of ``$k \le
  |t|$ implies the translation of the body'', with $W$ the width of $t$, and $\exists x \le |t|$ as
  the corresponding disjunction.
  \leanfile{Sagitta/Map/EspejoEF.lean}
\end{definition}

\begin{theorem}[Correctness of the translation, on the formulas of the witnessing calculus]
  \label{thm:arithmetic-trad-eval}
  \lean{Sagitta.Map.Extraccion.EspejoEF.trad_eval}
  \leanok
  \uses{def:arithmetic-translation, def:arithmetic-witness-calculus}
  Let $\sigma$ be an assignment of the propositional variables, $\varphi$ a formula of the class of
  Definition~\ref{def:arithmetic-translation}, $w$ a width, $E$ a vector of bits for each variable and
  $\rho$ an environment such that, for all $i$ and $j$, the value of $E(i)(j)$ under $\sigma$ is the
  bit $j$ of $\rho(i)$, and $\rho(i) < 2^w$ for all $i$. Then the translation $\mathrm{trad}(\varphi,
  w, E)$ is true under $\sigma$ if and only if $\varphi$ is true in the environment $\rho$.
  \source{\cite{Cook75, ParisWilkie85}.}
  \fidelity{Not checked against the printed source.}
  \leanfile{Sagitta/Map/EspejoEF.lean}
\end{theorem}

\begin{proof}
  \leanok
  Induction on $\varphi$. Atoms are handled by the correctness of the adder and the comparator, and the
  quantifiers by the correctness of the finite conjunction or disjunction over the values of the
  bound variable.
\end{proof}

\begin{theorem}[The induction instance for $N$ has a Frege derivation of $15N + 17$ lines]
  \label{thm:arithmetic-induction-frege}
  \lean{Sagitta.Map.Extraccion.EspejoInduccion.induccion_derivable}
  \leanok
  % uses-lean: Cslib.Logic.PL.Frege.DerivableInSize
  Let $V$ be a type of variables, $P : \mathbb{N} \to \mathrm{Proposition}(V)$ a family of
  propositional formulas and $N$ a number. Let the chain $\mathrm{cadena}(P, i, 0)$ be $\bot \to
  \bot$ and $\mathrm{cadena}(P, i, m + 1)$ be $(P_i \to P_{i+1}) \wedge \mathrm{cadena}(P, i + 1,
  m)$. The induction instance for $N$ is $(P_0 \wedge \mathrm{cadena}(P, 0, N)) \to P_N$. Then it has
  a derivation, without hypotheses, in the system of Frege of the machine library with at most $15N +
  17$ lines.
  \source{\cite{ParisWilkie85, Cook75}.}
  \fidelity{Not checked against the printed source.}
  \leanfile{Sagitta/Map/EspejoInduccion.lean}
\end{theorem}

\begin{proof}
  \leanok
  The derivation is built by hand under the antecedent, as a chain of $N$ steps of modus ponens; each
  step is $15$ lines: $6$ to take the step $P_k \to P_{k+1}$ out of the conjunction, $6$ to take out
  the rest, and $3$ for the modus ponens, with $17$ lines for the rest.
\end{proof}

\begin{definition}[The propositional formula that says that a circuit computes a function]
  \label{def:arithmetic-calcula}
  \lean{Sagitta.Map.Extraccion.TautologiaCircuitos.Calcula, Sagitta.Map.Extraccion.TautologiaCircuitos.Valida, Sagitta.Map.Extraccion.TautologiaCircuitos.Propaga, Sagitta.Map.Extraccion.TautologiaCircuitos.Coincide}
  \leanok
  \uses{def:arithmetic-translation}
  \own{}
  \infrastructure{}
  For $f : (\mathrm{Fin}\ n \to \{0,1\}) \to \{0,1\}$ and a number $s$, the propositional formula
  $\mathrm{Calcula}(f, s)$ is the conjunction of three parts. Its variables describe a circuit with
  $s$ gates over $n$ inputs: the type of each gate in two bits, its two wires, and the output wire,
  in registers of $n + s$ bits, and for each of the $2^n$ inputs the value of each gate. For an
  assignment $\sigma$ of the variables, $\mathrm{Valida}$ says that the output wire is less than $n +
  s$ and that each gate $k < s$ reads wires less than $n + k$; $\mathrm{Propaga}$ says that for every
  input $x < 2^n$ and gate $k < s$ the value of the wire $n + k$ on $x$ is the operation of the gate
  applied to the values of its two wires; $\mathrm{Coincide}$ says that for every $x < 2^n$ the value
  of the output wire on $x$ is $f(x)$.
  \leanfile{Sagitta/Map/TautologiaCircuitos.lean}
\end{definition}

\begin{theorem}[The formula of circuits, read assignment by assignment]
  \label{thm:arithmetic-evalb-calcula}
  \lean{Sagitta.Map.Extraccion.TautologiaCircuitos.evalB_calcula}
  \leanok
  \uses{def:arithmetic-calcula}
  Let $n$ and $s$ be numbers, $\sigma$ an assignment of the propositional variables and $f :
  (\mathrm{Fin}\ n \to \{0,1\}) \to \{0,1\}$. Then $\mathrm{Calcula}(f, s)$ is true under $\sigma$ if
  and only if $\mathrm{Valida}$, $\mathrm{Propaga}$ and $\mathrm{Coincide}$ hold for $\sigma$.
  \source{\cite{Razborov95}.}
  \fidelity{Not checked against the printed source.}
  \leanfile{Sagitta/Map/TautologiaCircuitos.lean}
\end{theorem}

\begin{proof}
  \leanok
  Unfold the conjunction, and read each of its three parts as the corresponding condition; the
  reading of the parts for propagation and for the output uses the validity.
\end{proof}

\section{The infinite length of $S^1_2$}
\label{sec:arithmetic-long}

This section is an exploration: it keeps what was proved about models of $S^1_2$ that are not
$\mathbb{N}$, and what was found not to follow. The statements about circuits are parameters, and the
second entry says that, with the way they are given, they have no content about circuits.

\begin{theorem}[A formula is not derivable from $S^1_2$ if and only if it fails at some element of a countable model]
  \label{thm:arithmetic-long-countermodel}
  \lean{Sagitta.Map.Acotada.LargoS12.no_deriva_iff_contraejemplo}
  \leanok
  \uses{def:arithmetic-proves, def:arithmetic-forall-exists, def:arithmetic-s12-model}
  Let $\varphi(x)$ be a formula of $L$ with one free variable. The closure of $\varphi$ is not
  derivable from $S^1_2$ if and only if there are an $L$-structure $M$ that is countable and a model
  of $S^1_2$, and an element $a \in M$ at which $\varphi(a)$ fails.
  \own{}
  \exploration{}
  \leanfile{Sagitta/Map/LargoInfinitoS12.lean}
\end{theorem}

\begin{proof}
  \leanok
  By the completeness theorem, a sentence is not derivable from $S^1_2$ if and only if it fails in a
  countable model of $S^1_2$, the term model; and the closure of $\varphi$ is true in $M$ if and only
  if $\varphi$ is true at every element.
\end{proof}

\begin{definition}[The instances of a formula, and what it means to express a bound]
  \label{def:arithmetic-expresa}
  \lean{Sagitta.Map.Acotada.LargoS12.instancia, Sagitta.Map.Acotada.LargoS12.Expresa, Sagitta.Map.Acotada.LargoS12.cotaInferiorSAT, Sagitta.Map.Acotada.LargoS12.cotaInferiorSATDesde}
  \leanok
  \own{}
  \exploration{}
  % uses-lean: Sagitta.Map.Tiempo.Finitud.sucNat, Sagitta.Map.CookS12.satBits
  For a formula $\varphi(x)$ of $L$ with one free variable and a number $n$, the instance
  $\varphi(\bar n)$ is the sentence obtained by putting for $x$ the numeral $\bar n = S(S(\cdots S(0)))$,
  with $n$ successors. For a formula $\varphi(x)$ and a predicate $Q$ on $\mathbb{N}$, $\varphi$
  \emph{expresses} $Q$ when for every $n$, $\varphi(n)$ is true in $\mathbb{N}$ if and only if $Q(n)$.
  Let $s(n)$ be the least number of gates of a Boolean circuit that decides SAT on the strings of
  length $n$. The lower bound $\mathrm{cotaInferiorSAT}(k, n)$ says $n^k < s(n)$, and
  $\mathrm{cotaInferiorSATDesde}(k, n_0, n)$ says that if $n_0 \le n$ then $\mathrm{cotaInferiorSAT}(k,
  n)$.
  \leanfile{Sagitta/Map/LargoInfinitoS12.lean}
\end{definition}

\begin{theorem}[Expressing a bound only asks that the bound be true]
  \label{thm:arithmetic-expresa-verdad}
  \lean{Sagitta.Map.Acotada.LargoS12.expresa_solo_pide_verdad}
  \leanok
  \uses{def:arithmetic-expresa, def:arithmetic-proves}
  Let $k$ and $n_0$ be numbers and $\varphi(x)$ a formula of $L$ with one free variable. Assume that
  for every $n$ the instance $\varphi(\bar n)$ is derivable from $S^1_2$. Then $\varphi$ expresses
  $\mathrm{cotaInferiorSATDesde}(k, n_0)$ if and only if $\mathrm{cotaInferiorSATDesde}(k, n_0, n)$
  holds for every $n$.
  \own{}
  \exploration{}
  \leanfile{Sagitta/Map/LargoInfinitoS12.lean}
\end{theorem}

\begin{proof}
  \leanok
  Derivable instances are true in $\mathbb{N}$, by soundness; so $\varphi(n)$ is true for every $n$, and
  both sides of the equivalence say that the bound holds for every $n$.
\end{proof}

\section{Proof systems}
\label{sec:arithmetic-proof-systems}

A proof system for a set of strings is defined as in the source; the entries below are two of its
elementary properties, and they do not involve any system in particular.

\begin{definition}[Proof systems, polynomially bounded, and polynomial simulation]
  \label{def:arithmetic-proof-system}
  \lean{Sagitta.Map.SistemaDePruebas.Sistema, Sagitta.Map.SistemaDePruebas.AcotadoPor, Sagitta.Map.SistemaDePruebas.PSimula}
  \leanok
  \source{\cite{CookReckhow79}.}
  \fidelity{Not checked against the printed source.}
  % uses-lean: Sagitta.Map.EsPolinomial
  A proof system for a set $L$ of lists of bits is a function $f$ from lists of bits to lists of bits
  that is computable in polynomial time by a Turing machine of Mathlib, such that $f(\pi) \in L$ for
  every $\pi$ and every $y \in L$ is $f(\pi)$ for some $\pi$. It is bounded by the polynomial $p$ when
  for every $y \in L$ there is $\pi$ with $f(\pi) = y$ and $|\pi| \le p(|y|)$. The system $S_2$
  polynomially simulates $S_1$ when there are a function $g$ computable in polynomial time and a
  polynomial $q$ such that $S_2.f(g(\pi)) = S_1.f(\pi)$ and $|g(\pi)| \le q(|\pi|)$ for every $\pi$.
  \leanfile{Sagitta/Map/SistemaDePruebas.lean}
\end{definition}

\begin{theorem}[A bounded proof system gives the form of $\cNP$]
  \label{thm:arithmetic-proof-np-form}
  \lean{Sagitta.Map.SistemaDePruebas.forma_np}
  \leanok
  \uses{def:arithmetic-proof-system}
  Let $S$ be a proof system for $L$ that is bounded by the polynomial $p$. Then for every list of bits
  $y$, $y \in L$ if and only if there is a list of bits $\pi$ with $|\pi| \le p(|y|)$ and $S.f(\pi) = y$.
  \source{\cite{CookReckhow79}.}
  \fidelity{Not checked against the printed source.}
  Only one direction of the equivalence between $\cNP$ and bounded proof systems is stated: the
  membership in $\cNP$ itself is not proved with the definitions of the repository.
  \leanfile{Sagitta/Map/SistemaDePruebas.lean}
\end{theorem}

\begin{proof}
  \leanok
  Forward, the bound gives the proof. Backward, $S.f(\pi) \in L$ by the correctness of the system.
\end{proof}

\begin{theorem}[Polynomial simulation is transitive]
  \label{thm:arithmetic-psimula-trans}
  \lean{Sagitta.Map.SistemaDePruebas.pSimula_trans}
  \leanok
  \uses{def:arithmetic-proof-system}
  Let $S_1$, $S_2$, $S_3$ be proof systems for $L$. If $S_3$ polynomially simulates $S_2$ and $S_2$
  polynomially simulates $S_1$, then $S_3$ polynomially simulates $S_1$.
  \source{\cite{CookReckhow79}.}
  \fidelity{Not checked against the printed source.}
  \leanfile{Sagitta/Map/SistemaDePruebas.lean}
\end{theorem}

\begin{proof}
  \leanok
  Compose the two translations, which is computable in polynomial time, and compose the polynomials
  that bound their lengths.
\end{proof}
